\documentclass[14 pt, fleqn, pstricks]{extarticle}

	\usepackage[frenchb]{babel}
	\usepackage[utf8]{inputenc}  
	\usepackage[T1]{fontenc}
	\usepackage{amssymb}
	\usepackage[mathscr]{euscript}
	\usepackage{stmaryrd}
	\usepackage{amsmath}
	\usepackage{tikz}
	\usepackage[all,cmtip]{xy}
	\usepackage{amsthm}
	\usepackage{varioref}
	\usepackage{geometry}
	\usepackage{tabularx}
	\geometry{a4paper}
	\usepackage{lmodern}
	\usepackage{hyperref}
	\usepackage{array}
	 \usepackage{fancyhdr}
	 \usepackage{pstricks,pst-plot,pst-tree,pstricks-add}
\usepackage{pst-eucl}% permet de faire des dessins de géométrie simplement
\usepackage{pst-text}
\usepackage{pst-node,pst-all}
\usepackage{pst-func,pst-math,pst-bspline,pst-3dplot}  %%% POUR LE BAC %%%
	 \usepackage{float}\usepackage{setspace}
\setlength{\mathindent}{1cm}
\renewcommand{\theenumi}{\alph{enumi})}
	\pagestyle{fancy}
	\theoremstyle{plain}
	\fancyfoot[C]{} 
	\fancyhead[L]{Contrôle }
	\fancyhead[R]{6 mai 2025}\geometry{
 a4paper,
 total={170mm,257mm}
 }
	
	
	\title{Exercices de calcul}
	\date{}
	\begin{document}
	 

\subsection*{Exercice 1 (1,5 x 2=3 points)}
Factorisez : 
\begin{enumerate}
\item $(x+1)^2 - 9$
\item $(2x+1)^2 - (x+2)^2$
\end{enumerate}

\subsection*{Exercice 2 (4x2=8 points) }	
Résoudre les équations suivantes :
\begin{enumerate}
\item $2x+10 = 16$
\item $8(x+5) = 44$
\item $5x + 3 = -2x + 7$
\item $(3x-1)(x+2) = 0$
\end{enumerate}

\subsection*{Exercice 3 (1+1+2+2=6 points) }	 
\textit{D'après l'exercice 4 du brevet Antilles-Guyane 2024.}\\
 On considère le programme de calcul ci-dessous :\\\fbox{\parbox{0.5\linewidth}{\setlength{\parskip}{.5cm}
\begin{itemize} 
	\item Choisir un nombre
	\item Mettre ce nombre au carré
	\item Soustraire le double du nombre de départ
	\item Soustraire 15
\end{itemize}
}}\newline

\begin{enumerate}
	\item Montrer que si on choisit $2$ comme nombre de départ, le résultat du programme est $-15$.
	\item On choisit $x$ comme nombre de départ. Exprimer le résultat du programme en fonction de $x$.
	\item Vérifier que l'on peut écrire ce résultat sous la forme $(x+3)(x-5)$.
	\item Déterminer les nombres à choisir au départ pour que le résultat du programme soit 0.
\end{enumerate}
     
     \subsection*{Exercice 4 : Mise en équation (2+2 points)}

Mettre en équation et résoudre les problème suivant : 
\begin{enumerate}
\item  Un père a 41 ans ; ses trois enfants sont âgés de 7, 9 et 13 ans. Dans combien d'années l'âge du
père sera-t-il égal à la somme des âges de ses enfants ?
\item La longueur d'un rectangle est supérieure à sa largeur de 21 m. Si on augmentait ses dimensions
de 4 m, la surface augmenterait de 492 m${}^2$.Trouver les dimensions de ce rectangle.
\end{enumerate}



 	\end{document}







